\documentclass[11pt,a4paper]{article}

\usepackage{graphicx}
\usepackage{xcolor}
\usepackage{geometry}
\usepackage{amsmath,amssymb,amsthm}
\usepackage{booktabs}
\usepackage{enumitem}
\usepackage{tabularx}
\usepackage{adjustbox}
\usepackage[numbers,sort&compress]{natbib}

% hyperref last among content packages
\usepackage[
    colorlinks=true,
    linkcolor=blue,
    citecolor=darkgray,
    urlcolor=blue,
    bookmarks=true,
    bookmarksnumbered=true,
    unicode=true
]{hyperref}

\geometry{a4paper, top=2.5cm, bottom=2.5cm, left=2.5cm, right=2.5cm}
\widowpenalty=10000
\clubpenalty=10000

\newtheorem{assumption}{Assumption}
\newcommand{\simplex}{\Delta_{K}}
\newcommand{\Rk}{R^{(k)}}
\newcommand{\Rho}{\rho}
\DeclareMathOperator*{\argmin}{arg\,min}

\begin{document}

\tableofcontents
\newpage

\section{Summary}\label{sec:summary}

This report derives and compares methods for estimating the ancestry
composition of a GWAS summary-statistics dataset from reference data alone.
Two families of estimators are considered. The first regresses the
per-variant effect-allele frequencies (EAF) reported in the summary
statistics on per-population reference frequencies; this channel is
implemented in the SMARTpred service and follows Summix
\citep{arriaga2021summix,stoneman2025summix2} and the bigsnpr ancestry
resources \citep{prive2018bigsnpr,privefl_vignette}. The second family,
which is the focus of this report, uses no allele frequencies at all: it
extracts the ancestry composition from the \emph{second moments of the
association statistics} --- products $z_i z_j$ and squared statistics
$\chi^2_i$ --- against per-ancestry LD reference panels. This is the channel
pioneered by DISTMIX \citep{lee2015distmix} and ZMIX
\citep{dennis2024zmix}, and it is grounded in the same moment identity that
underlies LD-score regression \citep{bulik2015ldsc}.

The practical motivation is threefold. First, some summary statistics ship
without a usable frequency column, or with one that cannot be trusted; the
LD channel needs only $\hat\beta_i$ and $\mathrm{SE}_i$. Second, the two
channels are statistically independent summaries --- the first moment of the
frequencies versus the second moment of the effects --- so their agreement
is itself a strong data-quality diagnostic, in the spirit of the
misspecification checks of \citet{prive2022identifying}. Third, for
cross-ancestry polygenic-score work such as the portability measurements of
this repository, an estimate of the ancestry composition of every input
dataset is a prerequisite for asserting that a ``target ancestry'' dataset
is what it claims to be.

The report is organised as follows. Section~\ref{sec:setting} fixes
notation. Section~\ref{sec:af} states the frequency-based estimator for
reference. Sections~\ref{sec:moments} and~\ref{sec:mixture} derive what the
association statistics know about the analysed population's LD.
Sections~\ref{sec:pairwise} and~\ref{sec:bilinear} define the two
LD-moment estimators. Sections~\ref{sec:uncertainty}
--~\ref{sec:failure} treat uncertainty, identifiability, and failure
modes; Section~\ref{sec:prior} positions the estimators against prior
work; Section~\ref{sec:implementation} maps them onto the SMARTpred and
ppb code bases.

\section{Setting and notation}\label{sec:setting}

Let $i = 1, \dots, m$ index variants (in practice a HapMap~3 panel
\citep{altshuler2010hapmap3} of roughly one million variants, harmonised
across references), and let $K$ reference ancestries be available with LD
correlation matrices $\Rk = (\Rho^{(k)}_{ij})$ on the standardized-genotype
scale, estimated from ancestry-matched reference panels (for example
1000~Genomes super-populations \citep{auton20151000g}; in the ldpred3
ecosystem, UK~Biobank references of the LDpred2 family
\citep{prive2020ldpred2}). The analysed GWAS dataset provides
\begin{equation}\label{eq:zdef}
  z_i \;=\; \frac{\hat\beta_i}{\mathrm{SE}_i},
  \qquad
  \chi^2_i \;=\; z_i^2,
\end{equation}
with effective sample size $n$ (per-variant $n_i$ where available), and
optionally a reported effect-allele frequency $f_i$. The unknown ancestry
composition is
\begin{equation}\label{eq:simplex}
  \pi \;=\; (\pi_1, \dots, \pi_K) \;\in\; \simplex
  \;=\; \Bigl\{ \pi : \pi_k \ge 0, \; \textstyle\sum_k \pi_k = 1 \Bigr\}.
\end{equation}

\section{The frequency channel (implemented)}\label{sec:af}

If the analysed population is a mixture of the reference populations with
little drift since admixture, its allele frequencies satisfy to first order
\begin{equation}\label{eq:afmix}
  f_i \;\approx\; \sum_{k=1}^{K} \pi_k \, p^{(k)}_i \;+\; \varepsilon_i,
\end{equation}
with $p^{(k)}_i$ the panel frequency of variant $i$'s counted allele in
ancestry $k$. The estimator minimises the residual sum of squares over the
simplex,
\begin{equation}\label{eq:afls}
  \hat\pi \;=\;
  \argmin_{\pi \in \simplex}
  \; \sum_{i=1}^{m} \Bigl( f_i - \sum_k \pi_k p^{(k)}_i \Bigr)^{2},
\end{equation}
after matching variants by identifier and allele pair (swapped alleles flip
$f_i$; strand-ambiguous variants are excluded, since frequency-based strand
resolution is unreliable exactly where ancestries differ). This is the
Summix estimator \citep{arriaga2021summix} up to optimisation details, and
it is what the SMARTpred ancestry mode implements with chromosome
jackknife standard errors and panel-confusability diagnostics. Its
requirement --- a trustworthy EAF column --- is precisely what the channels
of Sections~\ref{sec:pairwise}--\ref{sec:bilinear} dispense with.

\section{What the association statistics know}\label{sec:moments}

Write the marginal association statistic of variant $i$ as
\begin{equation}\label{eq:model}
  z_i \;=\; \frac{1}{\sqrt{n}}\, x_i^{\!\top} y,
  \qquad
  y \;=\; \sum_{j=1}^{m} x_j \gamma_j \;+\; \varepsilon,
\end{equation}
where $x$ is the standardized genotype matrix of the \emph{analysed}
population $A$ (mean zero, variance one per variant), the per-variant true
standardized effects satisfy $\operatorname{Var}(\gamma_j) = h^2/m$
(infinitesimal model), $\operatorname{Var}(\varepsilon) = 1 - h^2$, and $n$
is the effective sample size. Taking expectations over effects and noise,
with $R^{(A)} = (\Rho^{(A)}_{ij})$ the LD of the analysed population,
\begin{equation}\label{eq:pairmoment}
  \mathbb{E}\!\left[ z_i z_j \right]
  \;=\;
  \frac{n h^2}{m} \, \bigl( R^{(A)2} \bigr)_{ij}
  \;+\;
  \bigl( 1 - h^2 \bigr)\, \Rho^{(A)}_{ij}
  \qquad \text{for all } i, j,
\end{equation}
and, setting $i = j$ with $\Rho^{(A)}_{ii} = 1$ and
$\ell^{(A)}_i = \sum_j (\Rho^{(A)}_{ij})^2$ the LD score of variant $i$ in
the analysed population,
\begin{equation}\label{eq:ldscoreidentity}
  \mathbb{E}\!\left[ \chi^2_i \right]
  \;=\;
  1 \;-\; h^2 \;+\; \frac{n h^2}{m}\, \ell^{(A)}_i
  \;\approx\; a + b\, \ell^{(A)}_i ,
\end{equation}
the LD-score-regression identity \citep{bulik2015ldsc}: the intercept $a$
absorbs $1 - h^2$ plus any confounding inflation, and the slope
$b = n h^2 / m$ is the polygenic signal. Equations
\eqref{eq:pairmoment}--\eqref{eq:ldscoreidentity} state that the first two
moments of the association statistics are proportional to the LD structure
of the population that was analysed. That structure --- not the effect
sizes --- is what identifies the ancestry.

Two quantitative remarks. First, for a pair of variants in LD, both terms
of \eqref{eq:pairmoment} scale with the analysed population's own
correlation; for unlinked pairs both terms vanish, so unlinked pairs carry
no ancestry information. Second, the noise-covariance term dominates the
tagging term by roughly
\begin{equation}\label{eq:sigratio}
  \frac{(1 - h^2)}{(n h^2 / m)\,\bar\ell \cdot \Rho}
  \;\longleftrightarrow\;
  \frac{1 - h^2}{n h^2 \bar\ell / m}
  \;\approx\; \frac{0.75}{0.05}
  \quad \text{at } n = 10^5,\; h^2 = 0.25,\; m = 10^6,\; \bar\ell \approx 2:
\end{equation}
the pair-product moment is driven by correlated estimation noise, which is
present for \emph{any} trait with correct standard errors, including ones
with negligible polygenic signal. The $\chi^2$ moment, by contrast,
requires genuine signal ($h^2 > 0$), since at $h^2 = 0$ the expectation in
\eqref{eq:ldscoreidentity} is flat in $\ell$.

\section{Mixture model for the analysed population's LD}\label{sec:mixture}

Assume the analysed population is a random mating of the $K$ source
populations and that, for the variant separations used below, the haplotype
segment spanning a pair carries a single ancestry, drawn with probability
$\pi_k$ (admixture-LD is deferred to Section~\ref{sec:failure}). Then the
standardized-genotype correlation in $A$ is, to first order, the
ancestry-weighted mixture
\begin{equation}\label{eq:ldmix}
  \Rho^{(A)}_{ij} \;\approx\; \sum_{k=1}^{K} \pi_k\, \Rho^{(k)}_{ij}
  \qquad\text{for linked } (i,j).
\end{equation}
Substituting \eqref{eq:ldmix} into the dominant term of
\eqref{eq:pairmoment} keeps the pair moment \emph{linear in $\pi$}. The
$\chi^2$ route is different: the LD score squares the correlation, so
\begin{equation}\label{eq:quadmix}
  \ell^{(A)}_i
  \;=\; \sum_j \Bigl( \sum_k \pi_k \Rho^{(k)}_{ij} \Bigr)^{2}
  \;=\; \sum_{k} \sum_{k'} \pi_k \pi_{k'} \, \ell^{(kk')}_i,
  \qquad
  \ell^{(kk')}_i \;:=\; \sum_j \Rho^{(k)}_{ij} \Rho^{(k')}_{ij},
\end{equation}
which is linear in the \emph{products} $\pi_k \pi_{k'}$, not in $\pi$
itself. The bilinear scores $\ell^{(kk')}_i$ --- computable once from each
ordered pair of reference panels --- are the design variables of
Section~\ref{sec:bilinear}. Note that regressing $\chi^2$ on the plain
per-ancestry scores $\ell^{(k)}_i = \ell^{(kk)}_i$ instead is misspecified
whenever the cross-ancestry correlation vectors overlap, which they always
do; the cross terms $\pi_k \pi_{k'} \sum_j \Rho^{(k)}_{ij} \Rho^{(k')}_{ij}$
are exactly what that shortcut drops.

\section{Estimator A: pair products (linear in $\pi$)}\label{sec:pairwise}

For a set $P$ of sampled variant pairs with non-trivial reference LD, use
the dominant term of \eqref{eq:pairmoment}--\eqref{eq:ldmix} with an
unconstrained scale $s > 0$ absorbing $(1 - h^2)$ and the tagging term:
\begin{equation}\label{eq:paireq}
  \mathbb{E}\!\left[ z_i z_j \right]
  \;\approx\; s \sum_{k=1}^{K} \pi_k \Rho^{(k)}_{ij},
  \qquad (i,j) \in P .
\end{equation}
Estimate the composition by constrained least squares over the simplex,
profiling out the scale:
\begin{equation}\label{eq:pairls}
  (\hat\pi, \hat s)
  \;=\;
  \argmin_{\pi \in \simplex,\; s \ge 0}
  \; \sum_{(i,j) \in P}
  \Bigl( z_i z_j - s \sum_k \pi_k \Rho^{(k)}_{ij} \Bigr)^{2}.
\end{equation}
For fixed $\pi$, the optimal $s$ is the least-squares ratio
$\langle zz, L\pi \rangle / \lVert L\pi \rVert^2$ with $L$ the pair-by-ancestry
design matrix, so \eqref{eq:pairls} reduces to a one-dimensional search
over a simplex-constrained least-squares problem --- the same numerical
problem, in fact, that the frequency channel solves, and solvable by a
pure-NumPy active-set method. Requirements and properties:
\begin{enumerate}[itemsep=2pt]
  \item \textbf{No EAF column, no $h^2$ required.} The dominant noise term
        exists for any trait with correctly calibrated standard errors
        (Section~\ref{sec:moments}); what it needs is adequate $n$ and
        honest SEs.
  \item \textbf{Pairs must be linked.} Unlinked pairs have
        $\Rho^{(k)}_{ij} \approx 0$ for every $k$ and contribute only
        noise; take pairs within LD blocks, e.g.\ all pairs with
        $\max_k | \Rho^{(k)}_{ij} |$ above a floor, capped per block to
        bound cost.
  \item \textbf{Genome-wide $z$'s.} Because each pair is weakly
        informative, the estimator needs statistics across all chromosomes;
        a single region does not suffice.
\end{enumerate}
This is the estimator family of ZMIX \citep{dennis2024zmix}; DISTMIX
\citep{lee2015distmix} uses the same information indirectly, scoring each
ancestry by how well its LD imputes the observed $z$-pattern.

\section{Estimator B: bilinear LD-score regression}\label{sec:bilinear}

Regress the LDSC-style truncated statistic on the $K(K+1)/2$ bilinear
score vectors of \eqref{eq:quadmix}:
\begin{equation}\label{eq:bireg}
  \mathbb{E}\!\left[ \min(\chi^2_i,\, c) \right]
  \;=\;
  a \;+\; \sum_{k \le k'} b_{kk'} \, \ell^{(kk')}_i,
  \qquad c \approx 30,
\end{equation}
with non-negativity constraints $b_{kk'} \ge 0$ (they equal
$s\,\pi_k\pi_{k'}$ at truth) and an unconstrained intercept. Under
\eqref{eq:ldscoreidentity}--\eqref{eq:quadmix} the population values are
\begin{equation}\label{eq:bcoef}
  b_{kk'} \;=\; \frac{n h^2}{m}\; \pi_k \pi_{k'},
\end{equation}
so the composition is recovered from the row sums of the estimated
coefficient matrix, normalised:
\begin{equation}\label{eq:pirecover}
  \hat\pi_k
  \;=\;
  \frac{\sum_{k'} \hat b_{kk'}}
       {\sum_{k}\sum_{k'} \hat b_{kk'}} .
\end{equation}
The overall scale $n h^2/m$ cancels; $\hat\pi$ is then projected onto
$\simplex$ if needed, and the projection distance is itself a
specification diagnostic, since at truth $\hat B \approx \hat s\,
\hat\pi\hat\pi^{\top}$ is a scaled rank-one matrix. Per-variant sample
sizes enter multiplicatively in \eqref{eq:bcoef}: either restrict to
variants with comparable $n_i$ or fit with signal term $\ell_i / n_i$.
Estimator B needs genuine polygenic signal (flat $\chi^2$ at $h^2 = 0$
carries no LD profile) and well-powered statistics, but it is a single
per-variant regression over $m$ rows --- cheaper than the pair sampler of
Estimator A once the bilinear scores are precomputed.

\section{Uncertainty}\label{sec:uncertainty}

Residuals of both estimators are correlated across nearby variants and
pairs through LD, so i.i.d.\ standard-error formulas are wrong. Use the
delete-one-chromosome (or delete-one-LD-block) jackknife: refit on all but
one chromosome's variants (pairs), and form
\begin{equation}\label{eq:jackknife}
  \operatorname{SE}(\hat\pi_k)
  \;=\;
  \sqrt{ \frac{g-1}{g} \sum_{g'=1}^{g}
         \bigl( \hat\pi_k^{(-g')} - \bar{\hat\pi}_k \bigr)^{2} },
\end{equation}
with $g$ informative groups; the same construction the frequency channel
uses. For Estimator A, $\operatorname{Var}(z_i z_j) \approx
1 + \{\mathbb{E}[z_i z_j]\}^2$ under approximate normality suggests weights
$1/(1 + \hat\rho^2_{\text{fit},ij})$; the unweighted fit with jackknife is
the simpler default. For Estimator B the LDSC truncation already tames the
worst heteroskedasticity.

\section{Identifiability}\label{sec:identifiability}

In all three channels only the \emph{relative} proportions are identified:
the frequency channel is exact only up to the noise scale, and the LD
channels carry an unknown slope ($s$, or $n h^2/m$) that is not separable
from $\pi$ without external knowledge of $h^2$. This is the same target as
Summix and ZMIX --- a composition vector, not an heritability estimate.
The deeper identifiability constraint is geometric: the ancestry design
columns (frequencies $p^{(k)}$, pair correlations $\Rho^{(k)}_{ij}$, or
bilinear scores $\ell^{(kk')}$) can be severely collinear across
continents. LD patterns in particular decay with distance similarly
everywhere, so the LD design is worse conditioned than the frequency
design; the pairwise-column correlations and the Gram condition number
should be reported alongside any estimate (the confusability diagnostics
already implemented for the frequency panel generalise verbatim), and
near-collinear ancestries should be pooled rather than reported
separately.

\section{Assumptions and failure modes}\label{sec:failure}

\begin{assumption}[Random mating, segment-wise LD]\label{ass:mixture}
The mixture approximation \eqref{eq:ldmix} holds for linked pairs.
Admixture-LD violates it: the admixture process creates long-range
covariance of the form
$\pi_k (1 - \pi_k)\, d^{(k)}_i d^{(k)}_j$,
with $d^{(k)}_i$ the ancestry-informative frequency contrast, inflating
pair products at ancestry-informative markers and biasing $\hat\pi$. The
bias decays with admixture age; recent admixture is the hard case.
\end{assumption}

\begin{assumption}[Infinitesimal effects]\label{ass:infinitesimal}
$\operatorname{Var}(\gamma_j) = h^2/m$ with exchangeable effects on the
admixed-genotype scale. A trait driven by a few large loci violates the
moment structure of \eqref{eq:pairmoment}--\eqref{eq:bcoef}; the
truncation in \eqref{eq:bireg} mitigates this for Estimator B, less so for
Estimator A.
\end{assumption}

\begin{assumption}[Honest standard errors]\label{ass:se}
The dominant term of Estimator A is the correlated-noise covariance:
miscalibrated, winsorised, or meta-analysis-heterogeneity-distorted SEs
contaminate it directly. This is the channel's greatest sensitivity.
\end{assumption}

\begin{assumption}[No confounding beyond the intercept]\label{ass:confound}
Stratification and uncontrolled confounding add LD-score-shaped and
pair-product-shaped components (the LDSC intercept story of
\citet{bulik2015ldsc}). Free intercepts absorb only the mean part;
$\hat\pi$ drifts toward the LD profile of the confounding structure.
\end{assumption}

\begin{assumption}[Harmonised references]\label{ass:refs}
All $K$ references are built on the same variant panel with consistent
allele orientation, so pair correlations and bilinear scores are
comparable across ancestries. Panel construction from 1000~Genomes
\citep{auton20151000g} or UK~Biobank references \citep{prive2020ldpred2}
must pass through the same harmonisation as the target statistics.
\end{assumption}

\section{Relation to prior work}\label{sec:prior}

Table~\ref{tab:channels} summarises the channels.
\emph{Summix} \citep{arriaga2021summix} and its successor
\citep{stoneman2025summix2} deconvolve ancestry from allele frequencies
with a mixture model, and adjust frequencies to a target composition; the
bigsnpr ancestry resources \citep{prive2018bigsnpr, privefl_vignette}
implement the same idea through projection onto reference PCs followed by
simplex-constrained quadratic programming, with a worldwide 21-group
reference. \emph{DISTMIX} \citep{lee2015distmix} infers mixing proportions
in order to build ancestry-informed LD matrices for imputing summary
statistics; \emph{ZMIX} \citep{dennis2024zmix} regresses $z$-score products
on per-population genotype correlations by constrained least squares ---
Estimator A of Section~\ref{sec:pairwise} --- and reports, in simulation,
frequency-based methods near-perfect while the $z$-channel runs a few
percentage points off, worse for admixed targets: association statistics
simply carry less ancestry information than frequencies. The misspecification
checks of \citet{prive2022identifying} are the binary special case of all
of this --- ``does the reference match the statistics?'' --- where the
estimators here answer ``in what proportions?''. Estimator B's bilinear
scores are, to the best of our knowledge, new as an ancestry estimator; as
a construction they are the multi-ancestry generalisation of LD-score
regression \citep{bulik2015ldsc}, and nothing in this report claims
priority over methods we have not surveyed exhaustively.

\begin{table}[htbp]
\centering
\small
\caption{Ancestry-composition channels from summary statistics.}
\label{tab:channels}
\begin{tabular*}{\textwidth}{@{\extracolsep{\fill}} l l l l}
\toprule
Channel & Data used & Moment & Prior work \\
\midrule
Frequency & EAF column $f_i$ & first, \eqref{eq:afmix} &
  Summix \citep{arriaga2021summix}, \citep{stoneman2025summix2}, \\
 & & & bigsnpr \citep{privefl_vignette} \\
Pair products & $z_i z_j$, linked pairs & second, \eqref{eq:paireq} &
  ZMIX \citep{dennis2024zmix}, DISTMIX \citep{lee2015distmix} \\
Bilinear LD scores & $\chi^2_i$, all variants & second, \eqref{eq:bireg} &
  LDSC generalisation \citep{bulik2015ldsc}; new here \\
\bottomrule
\end{tabular*}
\end{table}

\section{Implementation mapping and recommendation}\label{sec:implementation}

In the SMARTpred service, the frequency channel is implemented
(\texttt{smartpred/ancestry.py}) with hash-pinned 1000~Genomes panels,
chromosome jackknife standard errors, and confusability reporting. Of the
LD channels, we recommend Estimator~A as the second channel to implement:
it needs no EAF column and no $h^2$, it reduces to the same
simplex-constrained least-squares solver already in place, and the
per-ancestry LD references and their LD-score sidecars already exist in
the service's registry. The pair set $P$ should be built from the
registered references' blocks (pairs within blocks with
$\max_k |\Rho^{(k)}_{ij}|$ above a floor, capped per block), with
delete-one-chromosome jackknife standard errors
\eqref{eq:jackknife}. Estimator~B is one further step once bilinear scores
$\ell^{(kk')}$ are precomputed per reference pair, and is the natural
benchmark arm against Estimator~A in a simulation study before either is
exposed on the public service; a block-wise Gaussian likelihood on the
full $z$-covariance remains research-grade (per-block cubic cost, plus
every misspecification of Assumptions~\ref{ass:mixture}--\ref{ass:confound}
at once).

For this repository, the estimators slot into the cross-ancestry
portability pipeline as a preflight check: estimate the composition of
every ``target ancestry'' dataset, verify it against the dataset's
declared ancestry, and refuse or flag measurements whose target is a
mixture the reference cannot separate --- the quantitative generalisation
of the reference-matching requirement stated by \citet{prive2022identifying}.

\section*{Acknowledgements}

The derivations in Sections~\ref{sec:moments}--\ref{sec:bilinear} were
worked out in discussion of the SMARTpred ancestry mode; the frequency
channel's implementation there served as the numerical template. All
citations were verified against the published records; any error in them
is ours.

\begin{thebibliography}{12}
\providecommand{\natexlab}[1]{#1}

\bibitem[Auton et~al.(2015)]{auton20151000g}
Auton, A., Brooks, L.~D., Durbin, R.~M., et~al. (1000~Genomes Project
Consortium) (2015).
A global reference for human genetic variation.
\emph{Nature}, 526(7571):68--74.

\bibitem[Altshuler et~al.(2010)]{altshuler2010hapmap3}
Altshuler, D.~M., Gibbs, R.~A., Peltonen, L., et~al. (International
HapMap~3 Consortium) (2010).
Integrating common and rare genetic variation in diverse human
populations.
\emph{Nature}, 467(7311):52--58.

\bibitem[Arriaga-MacKenzie et~al.(2021)]{arriaga2021summix}
Arriaga-MacKenzie, I.~S., Mazee, A.~M., Nguyen, A.~P., et~al., and
Hendricks, A.~E. (2021).
Summix: A method for detecting and adjusting for population structure in
genetic summary data.
\emph{American Journal of Human Genetics}.
\url{https://www.cell.com/ajhg/fulltext/S0002-9297(21)00221-4}

\bibitem[Bulik-Sullivan et~al.(2015)]{bulik2015ldsc}
Bulik-Sullivan, B.~K., Loh, P.-R., Finucane, H.~K., et~al. (2015).
LD Score regression distinguishes confounding from polygenicity in
genome-wide association studies.
\emph{Nature Genetics}, 47(3):291--295.

\bibitem[Dennis and Lee(2024)]{dennis2024zmix}
Dennis, T. and Lee, D. (2024).
ZMIX: estimating ancestry proportions using GWAS association Z-scores.
\emph{Bioinformatics Advances}, 4(1):vbae128.

\bibitem[Lee et~al.(2015)]{lee2015distmix}
Lee, D., Bigdeli, T.~B., Riley, B.~P., Fanous, A.~H., and
Bacanu, S.-A. (2015).
DISTMIX: direct imputation of summary statistics for unmeasured SNPs from
mixed ethnicity cohorts.
\emph{Bioinformatics}, 31(19):3099--3105.

\bibitem[Priv\'e et~al.(2018)]{prive2018bigsnpr}
Priv\'e, F., Luu, K., Blum, M.~G.~B., McGrath, J.~J., and
Visscher, P.~M. (2018).
Efficient analysis of large-scale genome-wide data with two R packages:
bigstatsr and bigsnpr.
\emph{Bioinformatics}, 34(16):2781--2787.

\bibitem[Priv\'e et~al.(2020)]{prive2020ldpred2}
Priv\'e, F., Arbel, J., Vilhj\'almsson, B.~J., et~al. (2020).
LDpred2: better, faster, and more accurate polygenic risk prediction.
\emph{American Journal of Human Genetics}, 107(6):1011--1021.

\bibitem[Priv\'e et~al.(2022)]{prive2022identifying}
Priv\'e, F., Arbel, J., Aschard, H., and Vilhj\'almsson, B.~J. (2022).
Identifying and correcting for misspecifications in GWAS summary
statistics and polygenic scores.
\emph{Human Genetics and Genomics Advances}, 3(4):100136.

\bibitem[Priv\'e(2022)]{privefl_vignette}
Priv\'e, F. (2022).
\emph{bigsnpr vignette: Ancestry --- estimating ancestry proportions from
allele frequencies} (software documentation).
\url{https://privefl.github.io/bigsnpr/articles/ancestry.html}

\bibitem[Stoneman et~al.(2025)]{stoneman2025summix2}
Stoneman, K.~P., Arriaga-MacKenzie, I.~S., Vasquez, K.~G., et~al., and
Hendricks, A.~E. (2025).
Characterizing substructure via mixture modeling in large-scale genetic
summary data.
\emph{American Journal of Human Genetics}.
\url{https://www.cell.com/ajhg/fulltext/S0002-9297(24)00449-X}

\end{thebibliography}

\hypersetup{
    pdftitle={Estimating Ancestry Composition from GWAS Summary Statistics},
    pdfauthor={Bjarni J. Vilhjalmsson},
    pdfsubject={Statistical genetics: ancestry estimation from summary statistics},
    pdfkeywords={ancestry composition, GWAS summary statistics, LD score, admixture, polygenic scores}
}

\end{document}
